\subsection*{Partie A}
	\begin{enumerate}
		\item On a :\quad
		$u_1 =  u_0 + 400 = \np{1025} +\np{1025} = \np{1425} $.

		\item Pour calculer un terme de la suite $(u_n)$ on ajoute 400 au terme précédent. Ainsi $(u_n)$ est une suite \textbf{arithmétique} de raison 400.

		\item \begin{tblr}[T]{colspec={*3{c}},vlines, hline{1,2,Z}={solid}}
			$n$ & $ u_n $ & calcul \\
			0 & \np{1025} & \\
			1 & \np{1425} & \np{1025} + 400\\
			2 & \np{1825} & \np{1425} + 400\\
			3 & \np{2225} & \np{1825} + 400\\
			4 & \np{2625} & \np{2225} + 400\\
			5 & \np{3025} & \np{2625} + 400\\
		\end{tblr}

		 Le dernier versement $u_5$ est une somme de \np[\euro{}]{3025}.
	\end{enumerate}

\subsection*{Partie B}

	\begin{enumerate}
		\item Augmenter de 10~\% c'est multiplier par $1,1$.

		$v_1 = 1,1\times \np{1550} = \np{1550} + 155 = \np{1705}$.

		\item Pour calculer un terme de la suite $(v_n)$ on multiplie le terme précédent par $1,1$.

		Ainsi $(v_n)$ est une suite \textbf{géométrique} de raison $1,1$.
	\end{enumerate}

\subsection*{Partie C}
	Pour la formule~1 les intérêts sont égaux à \quad $\np{12150.00} - \np{10000.00} = \np{2150.00}~\euro{}$.

	Pour la formule~2 les intérêts sont égaux à \quad $\np{11959.20} - \np{10000.00} = \np{1959.20}~\euro{}$.

	$\np{1929.20} < \np{2150}$. \quad Ainsi la formule~2 est plus avantageuse pour l'emprunteur car le montant des intérêts est moins important.
